\documentclass[11pt]{article}

\usepackage[a4paper,margin=1in]{geometry}
\usepackage{amsmath,amssymb,amsthm}
\usepackage{booktabs}
\usepackage{lmodern}
\usepackage[T1]{fontenc}
\usepackage{microtype}
\usepackage[numbers]{natbib}
\usepackage{hyperref}
\hypersetup{colorlinks=true,linkcolor=black,citecolor=black,urlcolor=blue}

\newtheorem{proposition}{Proposition}
\newtheorem{definition}{Definition}
\theoremstyle{remark}
\newtheorem*{remark}{Remark}

\newcommand{\TEmarg}{\mathrm{TE}_{\mathrm{marg}}}
\newcommand{\TEcrn}{\mathrm{TE}_{\mathrm{crn}}}
\newcommand{\DE}{\mathrm{DE}}
\newcommand{\ME}{\mathrm{ME}}
\newcommand{\Ex}{\mathbb{E}}
\newcommand{\TV}{\mathrm{TV}}

\title{Causal Attribution for Agentic Decisions:\\
Estimators, Coupling, and a Traceability Specification}
\author{Ajay Pravin Mahale}
\date{Draft of \today}

\begin{document}
\maketitle

\begin{abstract}
A provider of a high-risk AI system must keep records that make a decision
traceable, and for agentic systems it has not been established what those
records must contain for post-hoc causal attribution to be possible at all. We
give the estimator framework. We separate the marginal total effect that prior
work measures from a common-random-numbers total effect that isolates a step's
own contribution, add the natural direct effect under a pinned downstream, and
show the decomposition identity holds to machine precision on planted structure.
We derive the coupling that keeps the direct effect estimable once contexts
diverge, together with a closed form for its degradation, and we show that the
mediated share on which a natural ranking is built is not a share under
suppression: where the direct and mediated paths oppose it exceeds one and ranks
a suppressed component above a pure mediator, inverting the ordering it is meant
to induce. We pre-register the discrepancy experiment in full and publish the
pre-registration rather than a result, because the live regulated pipeline it
requires was not available in the study window; we state that absence here
rather than substituting a synthetic proxy for it. We contribute the
traceability specification such a filing would need, grounded in the verbatim
text of Articles 11 to 13 and Annex IV, whose obligations Regulation (EU)
2026/1744 deferred to 2 December 2027 on the Union's own record that the
apparatus for demonstrating conformity was not ready.
\end{abstract}

\section{Introduction}
\label{sec:intro}

\emph{To be written. Opens on the deferral and recital (40) of
Regulation (EU) 2026/1744 \citep{omnibus2026}, then the absence of an agreed
protocol for execution provenance \citep{survey2026}, then the contribution
list, with the differentiation from \citet{car2026} in this section rather than
in related work.}

\section{Setting and estimands}
\label{sec:setting}

\emph{To be written. Trajectory, node types, outcome $Y$, and the five-operator
intervention algebra adopted verbatim from \citet{car2026}.}

A trajectory is a sequence of steps $\tau = (v_1,\dots,v_T)$ terminating in a
decision $Y$. Step $k$ draws an action $a_k$ from a policy conditioned on the
prefix, using exogenous noise $u_k$. We write $a^{\mathrm{fact}}$ and
$u^{\mathrm{fact}}$ for the realised factual run, and every effect below is a
contrast against that run, as it must be: an effect is a contrast against
something that happened.

\section{Two total effects, and why the distinction is the crux}
\label{sec:two-te}

Fix a step $k$ and resample its action under the unchanged policy. Two estimands
follow, and they are not interchangeable.

\begin{definition}[Marginal total effect]
Intervene at $k$, then let every downstream step re-decide under \emph{fresh}
noise:
\[
  \TEmarg(k) \;=\; \Ex_{a'_k,\,u_{>k}}\!\left[\,Y \mid do(a_k := a'_k)\,\right]
              \;-\; y^{\mathrm{fact}} .
\]
\end{definition}

\begin{definition}[Common-random-numbers total effect]
Intervene at $k$, then let every downstream step re-decide under the
\emph{factual} noise $u_{>k} = u^{\mathrm{fact}}_{>k}$:
\[
  \TEcrn(k) \;=\; \Ex_{a'_k}\!\left[\,Y \mid do(a_k := a'_k),\,
                  U_{>k} = u^{\mathrm{fact}}_{>k}\,\right]
             \;-\; y^{\mathrm{fact}} .
\]
\end{definition}

The second is a counterfactual in Pearl's sense \citep{pearl2009}: same unit,
same exogenous noise, one variable changed. The first marginalises the noise
away and in doing so re-rolls every downstream decision. That is the confound
\citet{car2026} identifies and resolves with a point-of-commitment rule. We
resolve it by construction instead.

The construction is common random numbers, a standard variance-reduction
technique in simulation \citep{lawkelton2000}, and we claim no novelty in the
technique. Its use for rollout-based planning is recent \citep{yadav2026}. What
we claim is narrower and is the content of the next subsection: in this setting
the choice between marginalising the downstream noise and holding it fixed is
not a variance question at all. It decides whether a causally inert step is
distinguishable from a decisive one.

\subsection{A planted structure that separates them}

Consider a four-step chain with fidelity parameter $q$ and direct weight $w$:
\begin{align}
  a_0 &\sim \mathrm{Bern}(1/2) && \text{inert, no path to } y \nonumber\\
  a_1 &\sim \mathrm{Bern}(1/2) && \text{retrieval} \nonumber\\
  a_2 &\sim \mathrm{Bern}(1/2) && \text{inert, no path to } y \label{eq:scm}\\
  a_3 &= \begin{cases} a_1 & u_3 < q\\ 1-a_1 & \text{otherwise}\end{cases}
      && \text{executing tool call} \nonumber\\
  y   &= w\,a_1 + (1-w)\,a_3 . \nonumber
\end{align}
With $q = 0.9$, $w = 0.5$, and a factual run
$a^{\mathrm{fact}} = (0,1,1,1)$, $u_3^{\mathrm{fact}} = 0.770 < q$, so
$y^{\mathrm{fact}} = 1$. Table~\ref{tab:two-te} gives the closed forms, all
derived by hand before any estimator was run.

\begin{table}[t]
\centering
\begin{tabular}{llrrrr}
\toprule
step & planted role & $\TEmarg$ & $\TEcrn$ & $\DE$ & $\ME$ \\
\midrule
0 & inert, no path to $y$ & $-0.50$ & $0$      & $0$      & $0$ \\
1 & decisive retrieval    & $-0.50$ & $-0.50$  & $-0.25$  & $-0.25$ \\
2 & inert, no path to $y$ & $-0.05$ & $0$      & $0$      & $0$ \\
3 & executing step        & $-0.05$ & $-0.05$  & $-0.05$  & $0$ \\
\bottomrule
\end{tabular}
\caption{Analytic values on the planted chain \eqref{eq:scm}, $q=0.9$, $w=0.5$.
Under the marginal estimand the inert step 0 is numerically indistinguishable
from the decisive step 1, and the inert step 2 from the executing step 3.
Under CRN both inert steps are exactly zero.}
\label{tab:two-te}
\end{table}

The point of Table~\ref{tab:two-te} is the first column. Under marginal
run-forward, magnitude alone cannot separate a causally inert step from a
decisive one, which is precisely why a locus rule is needed on top of it. Under
CRN the inert steps return exactly zero and the planted structure is recovered
with no heuristic.

\begin{remark}
This is a statement about estimators on a known structure. It is not a statement
about any deployed system, and we do not write it as one.
\end{remark}

\section{The direct effect and the decomposition}
\label{sec:de}

\begin{definition}[Natural direct effect, pinned downstream]
Intervene at $k$ and pin every $j > k$ to the action it took in the factual run:
\[
  \DE(k) \;=\; \Ex_{a'_k}\!\left[\,Y \mid do(a_k := a'_k),\,
               a_{>k} := a^{\mathrm{fact}}_{>k}\,\right]
          \;-\; y^{\mathrm{fact}} .
\]
\end{definition}

This is the natural direct effect of \citet{pearl2001}, with the textbook
treatment in \citet[\S4.5.4--4.5.5]{pearl2009}; the concept predates the
counterfactual formalisation and priority belongs to \citet{robins1992}.

The mediated effect is the residual,
\begin{equation}
  \ME(k) \;=\; \TEcrn(k) - \DE(k),
  \label{eq:identity}
\end{equation}
which is an identity rather than an approximation. On the planted chain the
mediated share of step 1 is $|\ME|/|\TEcrn| = 0.25/0.5 = 0.5$, which equals
$1-w$ as it must, since the mediated path carries exactly weight $1-w$. That
identity is the closed-form check on the whole decomposition. Measured against
the derivation, the estimators agree within Monte Carlo error at every step and
the residual of \eqref{eq:identity} is $5.55 \times 10^{-17}$.

\subsection{Pin plausibility, reported rather than assumed away}

The direct-effect arm at step 1 pins $a_3 = 1$ while $a_1$ has been resampled.
When $a'_1 = 1$ the pin agrees with the policy with probability $q$; when
$a'_1 = 0$ it agrees with probability $1-q$. Averaging over
$a'_1 \sim \mathrm{Bern}(1/2)$ gives a pin plausibility of $(q + 1 - q)/2 = 0.5$
at $q = 0.9$. Where this quantity is small the direct-effect estimate rests on a
pinned continuation the policy would rarely have produced, and $\ME$ must be
read as bounded rather than as a point estimate. We report it alongside every
$\ME$.

\section{Coupling across divergent contexts}
\label{sec:coupling}

\citet[\S7]{car2026} states that a direct-effect arm needs common random numbers
across branches, calls this hard once contexts diverge, and leaves it as a
refinement. We give the closed form for what "hard" costs.

Token sampling is inverse-transform sampling from a categorical distribution. If
the uniform draw is a pure function of $(\text{run id}, \text{step},
\text{replicate})$ and never of the context, factual and counterfactual branches
consume identical draws. The noise is shared exactly. What is not shared exactly
is the benefit.

\subsection{Shared-\texorpdfstring{$u$}{u} coupling does not attain the bound}

A first implementation asserted that shared-$u$ inverse-transform sampling in
fixed index order attains the maximal-coupling bound $1 - \TV(p,q)$. That is
false, and the validation script rejected it at up to 215 Monte Carlo standard
errors before any of it reached this manuscript. Two branches agree at token $i$
exactly when $u$ lands in the intersection of the CDF intervals
$[P_{i-1},P_i)$ and $[Q_{i-1},Q_i)$, so
\begin{equation}
  A_{\mathrm{quant}}(p,q) \;=\;
    \sum_i \max\!\bigl(0,\; \min(P_i,Q_i) - \max(P_{i-1},Q_{i-1})\bigr)
  \;\le\; \sum_i \min(p_i,q_i) \;=\; 1 - \TV(p,q),
  \label{eq:quant}
\end{equation}
with equality only in degenerate cases. The mechanism behind the loss is that
mass moved between two tokens shifts every subsequent CDF boundary, so a
perturbation early in the index order decouples the entire tail.

Maximal coupling attains $1 - \TV(p,q)$, and that this is optimal is classical
\citep{levin2017}; we cite it rather than claim it. What we have not found
stated for this purpose, though it is an elementary computation, is the closed
form \eqref{eq:quant} for the shared-$u$ alternative and the size of the gap
between them at realistic divergence. Maximal coupling requires both $p$
and $q$ at draw time, so the factual branch's distribution must be carried
alongside the counterfactual one. That is one extra forward pass per step, and
it is available here because the factual run is being replayed anyway.

\begin{table}[t]
\centering
\small
\setlength{\tabcolsep}{5pt}
\begin{tabular}{lrrrrrr}
\toprule
logit shift & $\TV(p,q)$ & quant.\ bound & quant.\ emp. & max.\ bound &
max.\ emp. & prob-sorted \\
\midrule
0.00 & 0.0000 & 1.0000 & 1.0000 & 1.0000 & 1.0000 & 1.0000 \\
0.25 & 0.0856 & 0.7851 & 0.7839 & 0.9144 & 0.9142 & 0.5298 \\
0.50 & 0.1788 & 0.4121 & 0.4087 & 0.8212 & 0.8221 & 0.0858 \\
1.00 & 0.3969 & 0.2777 & 0.2740 & 0.6031 & 0.6075 & 0.0447 \\
2.00 & 0.4245 & 0.3057 & 0.3054 & 0.5755 & 0.5801 & 0.0446 \\
4.00 & 0.9369 & 0.0167 & 0.0177 & 0.0631 & 0.0654 & 0.0252 \\
\bottomrule
\end{tabular}
\caption{Vocabulary 32, 40{,}000 draws per point, $q$ formed by perturbing the
logits of $p$ with Gaussian noise of the stated scale. Every empirical value is
within 1.9 Monte Carlo standard errors of its closed form. The last column is
what a sampler that sorts the vocabulary by probability delivers, because
sorting breaks the fixed index order \eqref{eq:quant} depends on.}
\label{tab:coupling}
\end{table}

Table~\ref{tab:coupling} makes three design requirements binding. Use maximal
coupling, not quantile coupling: at $\TV = 0.179$ the gap is $0.821$ against
$0.412$. Fix the vocabulary order in the sampler: sorting by probability is the
default in common implementations and it collapses agreement to $0.086$ at that
same divergence, a factor of nine below what is achievable. And keep the uniform
draw a pure function of its key, since an advancing generator makes the draw
depend on call order, and call order differs between branches.

\subsection{What replay actually requires}

The construction above assumes the underlying forward pass is reproducible. It
is not reproducible by default. \citet{he2025} establish that the forward pass
is already run-to-run deterministic and that the usual explanation, concurrency
racing with floating-point non-associativity, is not the cause; the cause is
lack of batch invariance under varying server load. Their measurement is that
1000 completions at temperature zero from one prompt gave 80 unique
completions, identical for the first 102 tokens and diverging at the 103rd,
collapsing to one completion under batch-invariant kernels.

Three consequences bind on any replay harness. Temperature zero with a fixed
seed is not sufficient, and a harness that assumes otherwise will silently
mis-measure its own replay floor. The remedy is batch-invariant kernels rather
than serialisation. And determinism does not survive a change of GPU or
framework version, so a replay is valid only within one pinned inference stack,
which is the same discipline we apply to the analysis environment.

Where batch invariance cannot be obtained, $\ME$ inherits a noise floor that is
a property of the serving infrastructure rather than of the system under study,
and no effect smaller than the measured null-replay floor may be reported as an
effect.

\begin{remark}
This is a result about samplers, established on synthetic categorical
distributions. It has not been run against a language model.
\end{remark}

\section{Inconsistent mediation, and a pre-registered statistic built on it}
\label{sec:suppression}

A natural way to operationalise "how much of this step's influence travelled
through what it caused later steps to do" is the ratio $|\ME| / |\TEcrn|$. We
pre-registered a rank statistic on exactly that quantity. It is the wrong
quantity, and the mediation literature has said so for thirty years.

\subsection{The result is not ours}

In mediation analysis the ratio of the indirect effect to the total effect is
the \emph{proportion mediated}, and the configuration in which the direct and
indirect effects carry opposite signs is called \emph{inconsistent mediation},
with the mediator acting as a suppressor \citep{mackinnon2000,kenny_mediation}.
That the proportion mediated is unstable when the total effect is small, and can
exceed one or turn negative, is established in \citet{mackinnon1995}. We record
plainly that we did not derive this independently: we pre-registered a statistic
on the quantity, found the pathology on a constructed example, and located the
prior result on checking.

Restated in the notation of this paper, with no claim of novelty:

\begin{proposition}[Restatement of the proportion-mediated instability]
\label{prop:supp}
Let $\ME = \TEcrn - \DE$. If $\operatorname{sign}(\DE) \neq
\operatorname{sign}(\TEcrn)$, then
\[
  \frac{|\ME|}{|\TEcrn|} \;=\; \frac{|\TEcrn - \DE|}{|\TEcrn|}
  \;=\; 1 + \frac{|\DE|}{|\TEcrn|} \;>\; 1 ,
\]
unbounded above as $|\DE|$ grows in the opposing direction. A pure mediator has
$\DE = 0$ and attains exactly $1$, so a suppressed step outranks it.
\end{proposition}

\begin{proof}
$\ME$ is a difference, not a non-negative part, so opposite signs give
$|\TEcrn - \DE| = |\TEcrn| + |\DE|$. Divide by $|\TEcrn|$.
\end{proof}

\subsection{What is ours: the instrument, not the inequality}

Three things follow that the mediation literature had no occasion to state.

\paragraph{A pre-registered statistic ranked on a quantity that should not be
computed.} Our hypothesis H4 is a rank correlation on the mediated share.
Proposition~\ref{prop:supp} says that ranking places suppression above pure
mediation and mixes two opposite mechanisms into one score. The defect was in a
locked specification and was caught before any live data. We report it here
rather than in a footnote because the same construction is available to anyone
building attribution on a mediation decomposition, and the agent-attribution
literature has so far been built without reference to this methodology.

\paragraph{Inconsistent mediation has a concrete operational reading here.} In
an agent pipeline it is a retrieval that directly supports approval while
causing a later verification step to raise a flag. That is not a pathological
corner: it is a description of a component doing two opposing things, which is
common in a pipeline with a checking stage. Every structural model we used
before constructing the case had both paths pushing the same way, which is why
it had never been exercised.

\paragraph{The response is exclusion with a reported rate, not repair.} A share
is formed only where the decomposition splits a common-signed effect,
\[
  \operatorname{sign}(\DE) = \operatorname{sign}(\TEcrn)
  \quad\text{and}\quad
  |\DE| \le |\TEcrn| ,
\]
and is otherwise undefined and flagged. The rank statistic returns its exclusion
count alongside its value, so it cannot be reported without its exclusion rate.
Clamping $4/3$ to $1$ would repair an item that failed a validity check, which
our pre-registration forbids, and would merge suppression into pure mediation.
Steps with $|\TEcrn|$ below a pre-registered floor are undefined too, since a
share of no effect is undefined, but are not counted as suppression, so the
suppression rate is not inflated by inert steps.

\subsection{A worked instance}

Take the chain \eqref{eq:scm} with one sign flipped, $y = w\,a_1 - (1-w)\,a_3$,
at $w = 0.2$, $q = 0.9$. The same factual run gives
$y^{\mathrm{fact}} = -0.6$. Holding $u_3$ factual makes $a_3 = a'_1$, so
$y = (2w-1)a'_1$ and $\TEcrn(1) = +0.30$. Pinning $a_3 = 1$ gives
$\Ex[y] = -0.7$, so $\DE(1) = -0.10$, $\ME(1) = +0.40$ and the ratio is exactly
$4/3$. Measured over rollouts: $+0.3003$, $-0.1001$, $+0.4004$, ratio $1.3333$,
decomposition residual $5.55 \times 10^{-17}$. The estimators are correct; the
quantity built on them is not.

\begin{remark}[A previously computed value moved]
Applying the rule changed a rank correlation computed during pipeline
development from $\tau_b = +0.698$ to $+0.855$, with $80$ of $240$ steps
excluded, all inert and none suppressed. The earlier code assigned inert steps a
share of $0.0$, asserting their influence is entirely direct when they have no
influence at all. Those runs are synthetic smoke tests on planted generators
rather than results, and we report no rank correlation as a finding in this
paper.
\end{remark}

\section{The pre-registered experiment, published and not run}
\label{sec:prereg}

\emph{To be written from the pre-registration: H1 to H4, the joint
2-degree-of-freedom Wald primary contrast, the Plackett-Luce specification with
a decision-level cluster-robust sandwich, the reporting convention
$g = \beta/\lVert\beta\rVert_2$ and why a ratio to the causal coefficient is a
ratio-of-normals problem \citep{fieller1954}, the corpus rule, and the exclusion
rules. Stated plainly as not run here, and why.}

\section{An Annex IV traceability specification}
\label{sec:annexiv}

\emph{To be written from the verified regulatory basis.}

\section{Limitations}
\label{sec:limitations}

\emph{Short, because the principal limitation is in the abstract. Remaining: a
single planted generator family for validation, no live arm, no interaction
arm via the Shapley value, which is prior art \citep{car2026,castro2009}, and
coupling measured on synthetic categorical distributions rather than against a
language model.}

\section{Related work}
\label{sec:related}

\emph{To be written:
\citet{car2026,causalflow2026,agentracer2025,whoandwhen2025,ma2025,liao2026,survey2026,delegated2026,attriguard2026,causalarmor2026,civex2026,ebte2026}.}

\bibliographystyle{plainnat}
\bibliography{refs}

\end{document}
